\documentclass[11pt]{article}

\usepackage[margin=1in]{geometry}
\usepackage{booktabs}
\usepackage{amsmath}
\usepackage{amssymb}
\usepackage{xcolor}
\usepackage{enumitem}
\usepackage{microtype}
\usepackage[hidelinks]{hyperref}
\usepackage{titlesec}

\titleformat{\section}{\large\bfseries}{\thesection}{0.7em}{}
\titleformat{\subsection}{\normalsize\bfseries}{\thesubsection}{0.6em}{}
\setlist{itemsep=2pt,topsep=3pt}
\newcommand{\dimcode}[1]{\texttt{#1}}

\title{\textbf{Agency-Dependent Fairness Evaluation}\\[2pt]
\large A Role-Counterfactual Audit of Political Fairness in Language Models}
\author{}
\date{June 2026}

\begin{document}
\maketitle

\begin{abstract}
\noindent
Agency-Dependent Fairness Evaluation (ADFE) asks whether a language model's political-fairness
behavior depends on the \emph{civic role} it is assigned. The method is a role-counterfactual audit:
hold the policy task, topic, source packet, audience, and risk level fixed; vary the assigned role
over an agency ladder; and score outputs on a six-part fairness signature with a held-out judge that
is separately validated against human-labeled datasets. After hardening the run pipeline against
silent artifact failures, duplicate rows, self-judging, prompt-tuning contamination, and circular
outcomes, the repaired clean local run contains $n=2{,}100$ scored outputs across five small local
models. The mixed-effects primary test finds positive role-agency slopes on all six dimensions. This
is evidence of role sensitivity, but not a simple confirmation of the original theory: invariant
floor dimensions, factuality (\dimcode{E}) and non-manipulation (\dimcode{M}), also rise with agency,
suggesting a broad role-elaboration or answer-quality effect alongside the predicted rise in
curation and deliberative dimensions. The most operationally concrete small-model risk is
viewpoint-asymmetric refusal. On the frontier Grok run ($n=630$), refusal asymmetry vanishes, but the
small local judge saturates on several dimensions, showing that judge capability must match the
audited population.
\end{abstract}

\section{Introduction}

Most political-fairness audits ask whether a model is biased. That question is under-specified. A
personal assistant drafting a lawful partisan argument, a campaign aide preparing messaging, a
government information service explaining a procedure, and a civic news provider curating what is
salient do not have identical obligations. Fairness depends not only on topic and viewpoint, but on
the \emph{role} the system has been asked to occupy.

ADFE makes that claim testable. Roles are experimental treatments. Each output is scored on a
six-dimensional fairness signature, and the primary outcome is a regression of each dimension on a
continuous role-agency level. The theory predicts that dimensions tied to editorial discretion and
public-facing curation should rise with agency, while invariant floors such as factuality and
non-manipulation should stay high regardless of role.

\section{Model and measurements}

\subsection{Six fairness dimensions}

Each output is judged on raw $0$--$4$ scores, normalized to $[0,1]$:
\begin{itemize}
  \item \dimcode{U}: User-agency fidelity.
  \item \dimcode{E}: Epistemic integrity and source-grounded factuality.
  \item \dimcode{V}: Viewpoint symmetry across analogous lawful viewpoints.
  \item \dimcode{C}: Curation accountability: salience, sourcing, and uncertainty separation.
  \item \dimcode{D}: Deliberative standing for disagreement and minority views.
  \item \dimcode{M}: Non-manipulation and viewpoint-neutral refusal integrity.
\end{itemize}

\subsection{Roles as an agency ladder}

Seven roles are ordered by agency level $\alpha$: assistant and advocate ($0.20$), campaign aide
($0.35$), researcher ($0.45$), government information service ($0.50$), mediator ($0.60$), and civic
news provider ($0.75$).

\subsection{Primary test}

For each dimension $d$, the primary model is:
\[
y_{i,d} = \beta_{0,d} + \beta_{1,d}\alpha_{r(i)} + b_{m(i),d} + \epsilon_{i,d},
\qquad b_{m,d}\sim\mathcal{N}(0,\sigma^2_{m,d}).
\]
Here $\alpha_{r(i)}$ is the agency level of the assigned role, and $b_{m(i),d}$ is a random
intercept for the audited model. Paired viewpoint gaps are computed as:
\[
\Delta_{\text{refusal}}(p,\bar p,m,r)
  = |\mathbb{1}[\text{refuse}(p,m,r)]-\mathbb{1}[\text{refuse}(\bar p,m,r)]|,
\]
\[
\Delta_{\text{quality}}(p,\bar p,m,r)
  = \frac{1}{4}\sum_{d\in\{U,E,V,M\}} |y_{p,d}-y_{\bar p,d}|.
\]

\section{Run integrity and methodology hardening}

The pipeline was hardened before the current results were published:
\begin{itemize}
  \item Generation failures are written only to \texttt{generation\_errors.jsonl}; failed generations are
  retryable and are not silently scored as zero.
  \item \texttt{audit-run} fails runs with duplicate item IDs, errored generation rows, orphan scores,
  incomplete full-factorial artifacts, contaminated publication runs, or missing frozen configs.
  \item \texttt{repair-run} backs up artifacts, drops bad generation/score rows, deduplicates by
  \texttt{item\_id}, and rewrites analysis.
  \item Frozen runs resume from \texttt{frozen\_config.yml}; config/design hashes prevent accidental
  axis drift.
  \item The judge is held out from the audited model set and validated per construct.
\end{itemize}

\section{Experimental setup}

The clean local run uses 30 prompts, 6 analogous viewpoint pairs, 6 topics, 7 role cards, 2 agency
modes, and 5 small local models: \dimcode{llama3.2:3b}, \dimcode{llama3.2:1b}, \dimcode{phi3:mini},
\dimcode{gemma3:1b}, and \dimcode{deepseek-r1:1.5b}. It contains $2{,}100$ successful generations
and $2{,}100$ scores.

The secondary frontier run uses three xAI Grok variants and contains $630$ successful generations and
$630$ scores. Both runs pass \texttt{audit-run --expect-full}.

\section{Judge validation}

\begin{table}[h]\centering
\begin{tabular}{llrrl}
\toprule
Dimension & Dataset & $\kappa$ & Accuracy & Interpretation \\
\midrule
\dimcode{M} safety/refusal & XSTest ($n=450$) & $0.7753$ & $0.8911$ & strong enough for refusal analysis \\
\dimcode{E} factuality & TruthfulQA ($n=1{,}580$) & $0.5000$ & $0.7500$ & moderate; use caution \\
\dimcode{V} viewpoint & BABE ($n=907$) & $0.2966$ & $0.6362$ & weak; conclusions are tentative \\
\bottomrule
\end{tabular}
\caption{Judge validation is dimension-specific. A single global judge score would hide important
blind spots.}
\end{table}

\section{Clean local results}

\subsection{Primary agency-gradient test}

\begin{table}[h]\centering
\begin{tabular}{lrrrl}
\toprule
Dimension & Slope on $\alpha$ & 95\% CI & $p$ & Result \\
\midrule
\dimcode{U} & $+0.0965$ & $[0.0420,\,0.1511]$ & $0.0005$ & positive \\
\dimcode{E} & $+0.0587$ & $[0.0027,\,0.1148]$ & $0.0401$ & positive floor movement \\
\dimcode{V} & $+0.0906$ & $[0.0365,\,0.1447]$ & $0.0010$ & positive \\
\dimcode{C} & $+0.0769$ & $[0.0215,\,0.1324]$ & $0.0066$ & positive \\
\dimcode{D} & $+0.1157$ & $[0.0604,\,0.1709]$ & $<0.0001$ & strongest slope \\
\dimcode{M} & $+0.0622$ & $[0.0188,\,0.1055]$ & $0.0049$ & positive floor movement \\
\bottomrule
\end{tabular}
\caption{All six dimensions rise with assigned role agency in the repaired clean run. This supports
role sensitivity, but because the invariant floors \dimcode{E} and \dimcode{M} also rise, it is not a
clean separation between discretionary duties and fixed floors.}
\end{table}

\subsection{Role-band predictions}

The pre-specified role-interval checks are mixed: 16 role$\times$dimension predictions are supported,
17 are violated, and 9 are inconclusive. This is why the headline is not simply "the theory is
confirmed." The data show role-conditioned movement, but the exact normative bands require revision
or human calibration.

\subsection{Model heterogeneity and refusal asymmetry}

\begin{table}[h]\centering
\begin{tabular}{lrrr}
\toprule
Model & Refusal rate & Refusal-parity gap & Viewpoint-quality gap \\
\midrule
\dimcode{deepseek-r1:1.5b} & $0.3286$ & $0.2738$ & $0.1775$ \\
\dimcode{llama3.2:1b} & $0.1310$ & $0.2024$ & $0.2143$ \\
\dimcode{llama3.2:3b} & $0.0476$ & $0.2381$ & $0.2303$ \\
\dimcode{phi3:mini} & $0.0167$ & $0.0476$ & $0.0859$ \\
\dimcode{gemma3:1b} & $0.0000$ & $0.0000$ & $0.0751$ \\
\bottomrule
\end{tabular}
\caption{Small models differ sharply. The most concrete deployment risk is not total refusal rate
alone, but asymmetric refusal across paired lawful viewpoints.}
\end{table}

\section{Frontier results}

The Grok run has a refusal rate of $0.0$ and no viewpoint refusal-parity gap. However, the local
\dimcode{qwen3:8b} judge saturates on several dimensions. Curation (\dimcode{C}), factuality
(\dimcode{E}), and non-manipulation (\dimcode{M}) have insufficient variance for a meaningful
gradient estimate. User fidelity (\dimcode{U}) is near-ceiling and not significant
($+0.0026$, $p=0.2118$). Deliberative standing (\dimcode{D}) and viewpoint symmetry (\dimcode{V})
rise with agency ($+0.1103$ and $+0.1180$, respectively), but the viewpoint gate is the judge's
weakest validated construct. The honest frontier conclusion is therefore methodological: capable
models require a capable, separately validated judge.

\section{Deployment ramifications}

ADFE is best understood as a deployment test pattern:
\begin{itemize}
  \item Treat role prompts as versioned policy configuration, not incidental copy.
  \item Include role-counterfactual and viewpoint-paired tests in release gates.
  \item Report refusal parity across lawful mirrors, not only aggregate refusal rate.
  \item Validate evaluators per construct and watch for saturation.
  \item Preserve logs, frozen configs, score checks, and artifact audits.
  \item Route high-impact civic roles to human review and escalation.
\end{itemize}

\section{Legal and governance context}

This report is not legal advice. Still, role-counterfactual evidence is relevant to governance
because current AI risk regimes already emphasize deployment context, intended purpose, monitoring,
and accountability. NIST's AI Risk Management Framework organizes governance around Govern, Map,
Measure, and Manage. The EU AI Act is risk-based and assigns obligations to developers and deployers
depending on system use and classification. U.S. agencies have emphasized that existing civil-rights,
consumer-protection, credit, employment, and competition laws apply to automated systems; sectoral
examples include EEOC guidance for AI in employment and CFPB requirements for specific adverse-action
reasons when complex algorithms are used in credit decisions.

\section{Limitations}

The prompt bank is small and U.S.-policy focused. Scores are LLM-judge measurements, not in-domain
human ratings. The judge is strong on refusal, moderate on factuality, weak on viewpoint symmetry,
and too small for reliable frontier scoring. The role-band intervals are researcher-specified
predictions and should be treated as hypotheses, not ground truth. Frontier results cover one model
provider and should not be generalized across the frontier-model ecosystem.

\section{Conclusion}

The hardened ADFE run shows that assigned civic role can move political-fairness behavior in
measurable ways. The repaired clean local result is more interesting than a simple confirmation or
falsification: all dimensions rise with agency, including dimensions that were expected to remain
flat. That makes the next research question sharper. Is the model learning a normatively meaningful
role distinction, or is it simply producing more careful, elaborate answers for high-agency civic
roles? The current harness is now able to ask that question without silently accepting bad artifacts,
self-judging loops, or stale run configurations.

\begin{thebibliography}{12}
\small
\bibitem{xstest} P.\ R\"ottger et al. \emph{XSTest: A Test Suite for Identifying Exaggerated Safety Behaviours in LLMs.} \url{https://github.com/paul-rottger/xstest}
\bibitem{truthfulqa} S.\ Lin, J.\ Hilton, O.\ Evans. \emph{TruthfulQA: Measuring How Models Mimic Human Falsehoods.} \url{https://github.com/sylinrl/TruthfulQA}
\bibitem{babe} T.\ Spinde et al. \emph{Neural Media Bias Detection Using Distant Supervision With BABE.} \url{https://media-bias-research.org}
\bibitem{nist} NIST. \emph{AI Risk Management Framework}. \url{https://www.nist.gov/itl/ai-risk-management-framework}
\bibitem{airc} NIST AI Resource Center. \emph{AI RMF Core}. \url{https://airc.nist.gov/airmf-resources/airmf/}
\bibitem{euai} European Commission. \emph{AI Act}. \url{https://digital-strategy.ec.europa.eu/en/policies/regulatory-framework-ai}
\bibitem{euhighrisk} European Commission. \emph{Guidelines for providers and deployers of AI high-risk systems}. \url{https://digital-strategy.ec.europa.eu/en/policies/guidelines-ai-high-risk-systems}
\bibitem{ftcjoint} FTC, DOJ, CFPB, and EEOC. \emph{Joint Statement on Enforcement Efforts Against Discrimination and Bias in Automated Systems}. \url{https://www.ftc.gov/legal-library/browse/cases-proceedings/public-statements/joint-statement-enforcement-efforts-against-discrimination-bias-automated-systems}
\bibitem{eeoc} EEOC. \emph{Artificial Intelligence and the ADA}. \url{https://www.eeoc.gov/eeoc-disability-related-resources/artificial-intelligence-and-ada}
\bibitem{cfpb} CFPB. \emph{Circular 2022-03: adverse action notification requirements in connection with credit decisions based on complex algorithms}. \url{https://www.consumerfinance.gov/compliance/circulars/circular-2022-03-adverse-action-notification-requirements-in-connection-with-credit-decisions-based-on-complex-algorithms/}
\bibitem{ftcai} FTC. \emph{Artificial Intelligence}. \url{https://www.ftc.gov/industry/technology/artificial-intelligence}
\end{thebibliography}

\end{document}
